\origpage{19}{}
\hspace*{12pt}(a) $\hat{\Sigma}'$ \textit{is homeomorphic to} $\Sigma_0$\textit{, and}

\hspace*{12pt}(b) \textit{for} $v \in V(\Gamma')$, $v \in bd(\Sigma')$ \textit{if and only if either} $v \in V(A \cap B)$, \textit{or} $v \in V(\Gamma) \cap bd(\Sigma)$, \textit{or} $v \notin V(\Gamma)$ \textit{(that is,} $v$ \textit{is a new vertex produced by splitting).}

\textit{Then} $c(\Sigma') \ge c(\Sigma)$\textit{, and if equality holds then} $|V(\Gamma') \cap bd(\Sigma')| \ge p$.

\textit{Subproof.} As in (5), $\Gamma'$ has a tangle of order $\ge \theta - 2\theta_4$, and has a vital linkage with all its terminals in $bd(\Sigma')$ (from (b)). But
\begin{equation*}
|V(\Gamma') \cap bd(\Sigma')| \le |V(A \cap B)| + |V(\Gamma) \cap bd(\Sigma)| + 2\theta_4 \le p + 3\theta_4
\end{equation*}
and $\theta - 2\theta_4 \ge \theta_2(p + 3\theta_4)$, and so from (2) and (a), $c(\Sigma') \ge c(\Sigma)$. If equality holds then $\Sigma'$ is homeomorphic to $\Sigma$, and since $\theta - 2\theta_4 \ge \theta_3$ it follows that $|V(\Gamma') \cap bd(\Sigma')| \ge p$, from (3). This proves (6).

\hspace*{12pt}(7) \textit{There is no} $\Gamma$\textit{-normal proper line} $F \subseteq \Sigma$ \textit{with ends in different cuffs such that} $|F \cap V(\Gamma)| \le \theta_4$.

\textit{Subproof.} If there is such an $F$, let $A$ be null and $B = \Gamma$; let $\Sigma'$ be obtained from $\Sigma$ by cutting along $F$, and let $\Gamma'$ be obtained from $\Gamma$ by splitting appropriately the vertices of $\Gamma$ in $F$. (Since $\Gamma$ is loopless we can do so.) By (4.1)(i), $\hat{\Sigma}'$ is homeomorphic to $\Sigma_0$ and $c(\Sigma') = c(\Sigma) - 1$, contrary to (6). This proves (7).

Let $\mathcal{T}_1$ be the $\theta_4$-truncation of $\mathcal{T}$.

\hspace*{12pt}(8) $\mathcal{T}_1$ \textit{is respectful.}

\textit{Subproof.} Certainly $\Gamma$ is connected, by (4). Let $F \subseteq \hat{\Sigma}$ be a $\Gamma$-normal $O$-arc with $|F \cap V(\Gamma)| < \theta_4$. Suppose first that $F$ does not separate $\hat{\Sigma}$, let $\Sigma'$ be obtained from $\hat{\Sigma}$ by cutting along $F$, and let $\Gamma'$ be the drawing in $\Sigma'$ obtained from $\Gamma$ by splitting appropriately the vertices of $\Gamma$ in $F$. By (4.1)(iii) $\hat{\Sigma}'$ is simpler than $\Sigma_0$, contrary to (5). Thus $F$ separates $\hat{\Sigma}$. Let $\Sigma_1, \Sigma_2 \subseteq \hat{\Sigma}$ be surfaces such that $\Sigma_1 \cup \Sigma_2 = \hat{\Sigma}$ and $\Sigma_1 \cap \Sigma_2 = F$. Let $\Gamma_i = \Gamma \cap \Sigma_i$ ($i = 1, 2$). Since $(\Gamma_1, \Gamma_2)$ is a separation of $\Gamma$ of order $< \theta_4$, we may assume that $(\Gamma_1, \Gamma_2) \in \mathcal{T}_1 \subseteq \mathcal{T}$. By (5), $\hat{\Sigma}_2$ is not simpler than $\Sigma_0$, and so by (4.2)(iii), $\Sigma_1$ is a disc. This proves (8).
